\documentclass[11pt,letterpaper]{article}
\usepackage[T1]{fontenc}
\usepackage{lmodern,microtype,amsmath,amssymb,amsthm,booktabs}
\usepackage[margin=0.95in,headheight=14pt]{geometry}
\usepackage{fancyhdr,xcolor}
\usepackage[colorlinks=true,linkcolor=black,urlcolor=blue!55!black]{hyperref}
\newtheorem{theorem}{Theorem}
\newtheorem{lemma}[theorem]{Lemma}
\newtheorem{corollary}[theorem]{Corollary}
\newcommand{\oeis}[1]{\href{https://oeis.org/A#1}{A#1}}
\setlength{\parskip}{0.18em}
\raggedbottom

\hypersetup{pdftitle={Periods in the XOR triangle},pdfauthor={Blake Stone, with OpenAI Codex}}
\pagestyle{fancy}\fancyhf{}\lhead{\small Periods in the XOR triangle}\rhead{\small Blake Stone}\cfoot{\thepage}
\begin{document}
\thispagestyle{plain}
\begin{center}
{\LARGE\bfseries Least periods in the XOR Pascal triangle\par}
\vspace{0.5em}
Blake Stone\\
{\small Research manuscript prepared with OpenAI Codex}
\end{center}

\begin{abstract}
We determine the least period of every column and every active bit of the XOR triangle with diagonal \(1,2,3,\ldots\). A binomial formula also gives the complete rows at powers of two and the translated copies in the reflected triangle. These results prove the assertions of Franklin T. Adams-Watters in OEIS A195915 and A195916.
\end{abstract}
For \(1\le k\le n\), define \(T(n,n)=n\), use \(T(n,0)=0\), and set
\(T(n,k)=T(n-1,k-1)\mathbin{\oplus} T(n-1,k)\) for \(k<n\), where \(\oplus\) is bitwise exclusive or. The reflected triangle is \(U(n,k)=T(n,n+1-k)\).

\begin{theorem}\label{th:period}
In column \(k\), bit \(j\) (place value \(2^j\), \(j\ge0\)) is zero identically if \(2^j>k\). Otherwise it is purely periodic from the diagonal onward. Its least positive period is the smallest power of two at least \(k-2^j+1\). The whole column has least positive period \(2^{\lceil\log_2 k\rceil}\).
\end{theorem}
\begin{proof}
Over the field \(\mathbb F_2\), let \(d_k\) be bit \(j\) of the integer \(k\), and put
\[
 C_k(u)=\sum_{r\ge0} [\text{bit }j\text{ of }T(k+r,k)]u^r,\qquad C_0=0.
\]
The constant term is \(d_k\). For \(r\ge1\), the recurrence at offset \(r\) uses column \(k-1\) at offset \(r\) and column \(k\) at offset \(r-1\). Thus
\((1+u)C_k=C_{k-1}+d_k+d_{k-1}\). Incrementing \(k-1\) toggles bit \(j\) exactly when \(2^j\mid k\), giving
\[
 C_k(u)=\sum_{\substack{1\le s\le k\\ 2^j\mid s}}(1+u)^{-(k-s+1)}. \tag{1}
\]
An empty sum is zero. Otherwise the reduced denominator is \((1+u)^e\) with \(e=k-2^j+1\), since after multiplication by this denominator only the term \(s=2^j\) is nonzero at \(u=1\). The numerator has degree less than \(e\).

A purely periodic sequence of period \(D\) has generating function \(Q(u)/(1-u^D)\) with \(\deg Q<D\). Conversely any such fraction has period \(D\). Thus for (1) period \(D\) is equivalent to \((1+u)^e\mid1-u^D\); properness ensures the numerator degree condition after conversion. If \(D=2^v d\) with \(d\) odd then \(1-u^D=(1-u^d)^{2^v}\), and \(1-u^d\) has a simple root at one. Hence \(2^v\ge e\) is necessary and sufficient. The bit with \(j=0\) has \(e=k\), and all other bit periods divide its period.
\end{proof}

\begin{theorem}
For \(0\le j\) and \(n\ge k\ge1\), bit \(j\) of \(T(n,k)\) is
\[
 \sum_{\substack{1\le s\le k\\2^j\mid s}}\binom{n-s}{k-s}\pmod 2 .
\]
If \(n=2^J\), then \(T(n,k)=k\) at the powers of two \(k\le n\) and is zero elsewhere. The reflected triangle satisfies \(U(2^J+n,2^J+k)=U(n,k)\) whenever \(1\le k\le n\le2^J\).
\end{theorem}
\begin{proof}
Expand each term of (1) by the negative binomial series. At \(n=2^J\) use \(\binom{2^J-s}{2^J-k}\). In characteristic two the nonzero coefficients of \((1+x)^a=\prod_{i:\,a_i=1}(1+x^{2^i})\) occur precisely at indices whose set bits lie among those of \(a\). In \(J\) bits, \(2^J-s\) and \(2^J-k\) are the complements of \(s-1\) and \(k-1\). Thus the surviving \(s\) are exactly those for which \(s-1\) uses a subset of the set bits of \(k-1\). This already implies \(s\le k\).

To require \(2^j\mid s\), the lowest \(j\) bits of \(s-1\) must all be one. If any required bit is missing from \(k-1\), the count is zero. Otherwise each other set bit of \(k-1\) may be chosen freely, so the count is odd if and only if \(k-1=2^j-1\). This includes \(k=1,j=0\). Bits with \(2^j>k\) are already zero. The assertion about the row follows. Finally reflection sends both entries in the translation identity to the same column \(n+1-k\) of \(T\), with their row indices differing by \(2^J\). Theorem~\ref{th:period} shows that this is a period.
\end{proof}

For example the fourth column, starting with \(T(4,4)\), repeats \(4,4,6,7\) as determined by (1).
The original statements, including the claimed period and the reflected copies, are recorded by F. T. Adams-Watters in \oeis{195916} and \oeis{195915} (2011).

\paragraph{Independent formal proof.}
The accompanying Lean file proves the exact column and bit periods, the reflected translations, and the rows at every power of two from the defining recurrence. Its first identity is
\[
 T(n+2^h,k)=T(n,k)\mathbin{\oplus}T(n,\max(0,k-2^h))\qquad(k\le n).
\]
Applying a shift twice cancels the two middle terms. If \(2^h<k\le2^{h+1}\), the shift by \(2^h\) already fails at the diagonal because \(T(k,k-2^h)=k-2^h\ne0\). For dyadic rows the formal proof also gives \(T(2^{h+1},k)=T(2^h,k)\) away from the new diagonal, extending rows by zeros on the right.
\end{document}
